%==============================================================================
% CHAPTER 3: Spectral Methods in Microeconomics
%==============================================================================
\chapter{Spectral Methods in Microeconomics}
\label{ch:spectral} \cite{golub2026spectral}

\begin{flushright}
\textit{Benjamin Golub}\\
Communicated by \textit{Notices} Associate Editor Daniela De Silva
\end{flushright}

\epigraph{``Matrices often appear in formal models of social and economic behavior... Spectral theory offers powerful tools for understanding matrices, and economic modelers have leveraged these tools to gain considerable insight.''}{--- Benjamin Golub}

\begin{abstract}
Eigenvalues and eigenvectors are the hidden skeleton of networked systems. This chapter shows how Perron-Frobenius \index{Perron-Frobenius theorem}theory, Katz-Bonacich centrality, \index{centrality measures}and spectral decomposition unify opinion dynamics, network games, and market equilibrium. The same mathematics describes how beliefs spread, how firms compete, and how central banks should act, \cite{golub2026spectral}.
\end{abstract}

%==============================================================================
\section{Introduction: Why Spectral Methods?}

\noindent\textit{Mathematical notation follows standard conventions. Key terms are defined in the glossary.}\n
\label{sec:spec:intro}

In microeconomics, agents interact: consumers learn from neighbors, firms compete on networks, banks lend to each other. These interactions are encoded in \textbf{matrices}. The spectral properties of these matrices---especially the dominant eigenvalue and eigenvector---determine the system's long-run behavior, \cite{golub2026spectral}.

Consider a row-stochastic matrix $\bm{W}$ (nonnegative, rows sum to 1) representing social influence. The DeGroot \index{DeGroot model}model of opinion dynamics:
\[
\bm{x}(t+1) = \bm{W} \bm{x}(t).
\]
If $\bm{W}$ is irreducible and aperiodic, $\bm{x}(t) \to \bm{\pi}^\top \bm{x}(0) \bm{1}$ where $\bm{\pi}$ is the \textbf{left Perron vector} (stationary distribution). Convergence rate is governed by $|\lambda_2|$, the second-largest eigenvalue modulus.

This single fact unifies:
\begin{itemize}
    \item \textbf{Opinion dynamics}: How fast do beliefs converge?
    \item \textbf{Centrality}: Who is influential? (Katz-Bonacich = Perron vector)
    \item \textbf{Network games}: Equilibrium actions = $(\bm{I} - \beta \bm{G})^{-1} \bm{a}$; spectral radius of $\beta \bm{G}$ determines uniqueness, \cite{golub2026spectral}.
    \item \textbf{Information aggregation}: ``Wisdom of crowds'' iff $\max_i \pi_i \to 0$ as $n \to \infty$.
\end{itemize}

%==============================================================================
\section{Perron-Frobenius Theory for Economists} \cite{golub2026spectral}
\label{sec:spec:pf}

\begin{theorem}[Perron-Frobenius (Irreducible Nonnegative)] \cite{golub2026spectral}
Let $\bm{A} \in \mathbb{R}^{n \times n}$, $a_{ij} \ge 0$, irreducible. Then:
\begin{enumerate}
    \item $\rho(\bm{A}) > 0$ is a simple eigenvalue (the \textbf{Perron root}).
    \item There exist $\bm{u} > 0$, $\bm{v} > 0$ (right/left Perron vectors) with $\bm{A}\bm{u} = \rho\bm{u}$, $\bm{v}^\top\bm{A} = \rho\bm{v}^\top$.
    \item All other eigenvalues satisfy $|\lambda| < \rho$.
    \item For any $\bm{x} > 0$, $\bm{A}^t \bm{x} / \rho^t \to \bm{u} (\bm{v}^\top \bm{x} / \bm{v}^\top \bm{u})$.
\end{enumerate}
\end{theorem}

\begin{definition}[Irreducibility]
$\bm{A}$ is \textbf{irreducible} iff its directed graph is strongly connected: for any $i,j$, there is a path $i \to j$. Equivalently, no permutation makes $\bm{A}$ block upper-triangular.
\end{definition}

In economics, irreducibility = ``the network is connected.'' If the economy has isolated sectors, spectral theory applies block by block, \cite{golub2026spectral}.

%==============================================================================
\section{Application 1: DeGroot Learning and Social Influence}
\label{sec:spec:degroot}

Agents update beliefs by averaging neighbors:
\[
x_i(t+1) = \sum_j w_{ij} x_j(t), \quad \bm{W} \text{ row-stochastic}.
\]

\begin{proposition}
If $\bm{W}$ is irreducible and aperiodic, $\bm{x}(t) \to \bm{\pi}^\top \bm{x}(0) \bm{1}$ where $\bm{\pi}$ is the unique stationary distribution ($\bm{\pi}^\top \bm{W} = \bm{\pi}^\top$, $\bm{\pi} > 0$, $\sum \pi_i = 1$).
\end{proposition}

The \textbf{social influence} of agent $i$ is $\pi_i$. This is exactly \textbf{eigenvector centrality} (left Perron vector of $\bm{W}$).

\begin{example}[Star network]
Agent 1 at center, all others peripheral. $\bm{W}_{1j} = 1/(n-1)$, $\bm{W}_{i1} = 1$ for $i>1$. Then $\pi_1 = 1/2$, $\pi_i = 1/(2n-2)$ for $i>1$. The center has influence $\sim 1/2$ regardless of $n$!
\end{example}

\begin{definition}[Wisdom of Crowds]
A sequence of networks is \textbf{wise} if $\max_i \pi_i^{(n)} \to 0$. Golub-Jackson (2010): wisdom holds iff no agent has ``disproportionate influence.''
\end{definition}

%==============================================================================
\section{Application 2: Network Games and Bonacich Centrality}
\label{sec:spec:games}

Agents choose actions $a_i \in \mathbb{R}$. Payoff:
\[
u_i(a_i, \bm{a}_{-i}) = a_i \left( \alpha_i - \frac{1}{2} a_i + \beta \sum_j g_{ij} a_j \right).
\]
Best response: $a_i = \alpha_i + \beta \sum_j g_{ij} a_j$. In vector form:
\[
\bm{a} = \bm{\alpha} + \beta \bm{G} \bm{a} \quad \Rightarrow \quad \bm{a}^* = (\bm{I} - \beta \bm{G})^{-1} \bm{\alpha},
\]
provided $\rho(\beta \bm{G}) < 1$.

\begin{definition}[Katz-Bonacich Centrality]
For $\bm{G} \ge 0$ and $\beta < 1/\rho(\bm{G})$,
\[
\bm{b}(\beta, \bm{G}) = (\bm{I} - \beta \bm{G})^{-1} \bm{1} = \sum_{k=0}^\infty \beta^k \bm{G}^k \bm{1}.
\]
Entry $b_i$ counts weighted walks starting at $i$; longer walks discounted by $\beta^k$.
\end{definition}

The equilibrium action is $\bm{a}^* = \bm{b}(\beta, \bm{G}) \odot \bm{\alpha}$ (if $\bm{\alpha}$ is scalar). \textbf{Equilibrium actions are proportional to centrality.}

%==============================================================================
\section{Application 3: Wisdom, Centralization, and Policy}
\label{sec:spec:policy}

The Golub-Jackson \textbf{prominent sequence} concept: a family of networks $(\bm{W}^{(n)})$ with $n$ agents. They are wise iff the most influential agent's weight $\to 0$.

\begin{theorem}[Wisdom Criterion]
$(\bm{W}^{(n)})$ is wise iff $\max_i \pi_i^{(n)} \to 0$. For symmetric $\bm{W}$, this is equivalent to $\lambda_2(\bm{W}) \to 1$ (the spectral gap closes), \cite{golub2026spectral}.
\end{theorem}

Policy insight: if a regulator wants to improve information aggregation, they should \textbf{reduce the spectral gap}---i.e., make the network less centralized. This contradicts the intuition that ``more connectivity is always better.'' A star network is highly connected but unwise; a line network is poorly connected but wise, \cite{golub2026spectral}.

%==============================================================================
\section{Code: Spectral Analysis of Networks}
\label{sec:spec:code}

\begin{lstlisting}[style=julia,caption=Network spectral analysis] \cite{golub2026spectral}
using LinearAlgebra, SparseArrays, LightGraphs
using Statistics

function analyze_network(W; tol=1e-12)
    n = size(W, 1)
    
    # Check irreducibility
    G = DiGraph(W .> tol)
    is_irred = is_strongly_connected(G)
    
    # Perron-Frobenius (left eigenvector for row-stochastic) \cite{golub2026spectral}
    eig = eigen(W')
    idx = argmax(abs.(eig.values))
    lambda1 = real(eig.values[idx])
    pi = real(eig.vectors[:, idx])
    pi = pi / sum(pi)  # normalize
    
    # Second eigenvalue modulus
    lambda2 = maximum(abs.(eig.values[setdiff(1:end, idx)]))
    spectral_gap = lambda1 - lambda2 \cite{golub2026spectral}
    
    # Katz-Bonacich centrality for beta < 1/rho(G)
    # Assuming G is the adjacency (not stochastic)
    function katz_centrality(G, beta)
        rho = maximum(abs.(eigvals(G)))
        if beta >= 1/rho
            error("beta too large: beta*rho = $(beta*rho)")
        end
        (I - beta*G) \ ones(n)
    end
    
    # Wisdom metrics
    max_influence = maximum(pi)
    herfindahl = sum(pi.^2)  # concentration index
    
    return (
        lambda1 = lambda1,
        lambda2 = lambda2,
        spectral_gap = spectral_gap, \cite{golub2026spectral}
        pi = pi,
        max_influence = max_influence,
        herfindahl = herfindahl,
        is_irred = is_irred
    )
end

# Example: Star network
n = 10
W = zeros(n,n)
W[1, 2:end] .= 1/(n-1)
for i = 2:n
    W[i, 1] = 1.0
end

res = analyze_network(W)
println("Star network: max_influence = ", res.max_influence)
println("Spectral gap = ", res.spectral_gap) \cite{golub2026spectral}

# Example: Line network
W2 = zeros(n,n)
for i = 1:n
    neighbors = []
    i > 1 && push!(neighbors, i-1)
    i < n && push!(neighbors, i+1)
    W2[i, neighbors] .= 1/length(neighbors)
end
res2 = analyze_network(W2)
println("Line network: max_influence = ", res2.max_influence)
println("Spectral gap = ", res2.spectral_gap) \cite{golub2026spectral}
\end{lstlisting}

%==============================================================================
\section{Bridge to Chapter 4}
\label{sec:spec:bridge}

Spectral methods analyze \textbf{linear} network dynamics: $\bm{x}(t+1) = \bm{W}\bm{x}(t)$. Chapter 4 introduces \textbf{nonlinear} network dynamics: deep neural networks. Remarkably, ReLU networks are \emph{piecewise linear}---on each linear region, they are exactly a matrix multiplication. The geometry of deep learning is the geometry of how these linear regions tessellate the input space.

%==============================================================================
\section{Further Reading}
\label{sec:spec:refs}

\begin{itemize}
    \item Golub, B. (2025). \textit{Spectral Methods in Microeconomics}. Notices of the AMS, 72(6).
    \item Golub, B., Jackson, M.O. (2010). \textit{Naive Learning in Social Networks and the Wisdom of Crowds}. American Economic Journal: Microeconomics.
    \item Ballester, C., Calvó-Armengol, A., Zenou, Y. (2006). \textit{Who's Who in Networks. Wanted: The Key Player}. Econometrica.
    \item Jackson, M.O. (2010). \textit{Social and Economic Networks}. Princeton University Press.
    \item Bonacich, P. (1987). \textit{Power and Centrality: A Family of Measures}. American Journal of Sociology.
\end{itemize}

%==============================================================================
\section{Exercises}
%==============================================================================
% Solutions
%==============================================================================
\section{Solutions}
\label{sec:spec:solutions}

\begin{solution}
\textbf{Exercise 1 (Star network):} For the star graph with center node 1 and $n-1$ leaf nodes, the adjacency matrix has eigenvalues $\sqrt{n-1}$, $-\sqrt{n-1}$, and $0$ with multiplicity $n-2$. The stationary distribution is $\pi_1 = 1/2$ and $\pi_i = 1/(2(n-1))$ for $i \neq 1$.
\end{solution}

\begin{solution}
\textbf{Exercise 2 (Katz-Bonacich):} The Katz-Bonacich centrality is $b = (I - \beta G)^{-1}\mathbf{1}$. In the DeGroot model with $G = W$ (row-stochastic), this gives $b = (I - \beta W)^{-1}\mathbf{1} = \sum_{k=0}^{\infty} \beta^k W^k \mathbf{1}$. Each term $\beta^k W^k \mathbf{1}$ represents $k$-step influence weighted by $\beta^k$.
\end{solution}

\begin{solution}
\textbf{Exercise 3 (Wisdom test):} For $p_n = \log n / n$, the graph is connected with high probability, and the average degree is $\log n \to \infty$. The wisdom test holds: $\max_i \pi_i \to 0$. For $p_n = 1/\sqrt{n}$, the graph is disconnected with high probability, and the wisdom test fails.
\end{solution}

\begin{solution}
\textbf{Exercise 4 (Heterogeneous $\beta_i$):} For heterogeneous interaction strengths, the equilibrium is $x^* = (I - \text{diag}(\beta) G)^{-1} a$. A unique equilibrium exists when $\rho(\text{diag}(\beta) G) < 1$. This condition can be checked using the spectral radius.
\end{solution}

\label{sec:spec:exercises}

\begin{exercise}
For the star network with $n$ agents, compute the exact stationary distribution $\pi$ and verify $\pi_1 = 1/2$ for all $n \ge 3$.
\end{exercise}

\begin{exercise}
Show that for a row-stochastic matrix $\bm{W}$, the Katz-Bonacich centrality with $\bm{G} = \bm{W}$ and $\beta < 1$ is $\bm{b} = (\bm{I} - \beta\bm{W})^{-1}\bm{1}$. What does this represent in the DeGroot model?
\end{exercise}

\begin{exercise}
Implement the Golub-Jackson wisdom test: generate a sequence of Erdős-Rényi graphs $G(n, p_n)$ with $p_n = \log n / n$ and $p_n = 1/\sqrt{n}$. Which sequence is wise?
\end{exercise}

\begin{exercise}[Research]
Extend the network game to heterogeneous $\beta_i$ (different interaction strengths). When does a unique equilibrium exist?
\end{exercise}